As dataset sizes in AI/HPC workloads increase, data compression could become
useful to speed up overall execution time by reducing the time taken to transfer
between nodes on the network. However, this is not useful if the time
\emph{cost} of compression/decompression outweighs the time \emph{savings} in
data transmission over the network. This paper explores methods of compressing
data --- using OpenZL, a recent lossless compression framework --- in various
scenarios, i.e., compression before execution, compression during execution, and
no compression. I compare and analyze the time taken with various compression
algorithms and datasets. With the scenarios and datasets tested, \{all
scenarios; some scenarios; no scenarios\} exhibited speedup as a result of
compression with OpenZL.
